%% kapa: exam=1 round=25 year=2014 period=1 form=A subject=civil_law subject_ko=민법 seq=26 number=26
%% source: Q-Net 감정평가사 기출문제 (cst00309) artlSeq=5101207 file=extracted/1008/5101207/2014년 제25회 감정평가사 1차시험 기출문제/1교시(A형).hwp
%% conversion: kordoc 4.24.51 -> tools/gen1.py (auto); review=proofread-sonnet; flags=none
\begin{question}{26}
소멸시효의 중단과 정지에 관한 설명으로 옳지 않은 것은? (다툼이 있으면 판례에 의함)
\begin{choices}
  \item 승인은 시효의 진행이 개시된 후에만 할 수 있고, 그 전에 승인하더라도 시효가 중단되지 않는다.
  \item 여러 차례의 최고가 있은 후, 재판상 청구가 있더라도 그 시효중단의 효력은 항상 최초의 최고를 한 때에 발생한다.
  \item 가압류에 의한 시효중단의 효력은 가압류의 집행보전의 효력이 존속하는 동안 계속된다.
  \item 재판상의 청구로 중단된 시효는 재판이 확정된 때로부터 새로이 진행한다.
  \item 천재 기타 사변으로 인하여 시효를 중단할 수 없을 때에는, 그 사유가 종료한 때로부터 1개월 내에는 시효가 완성하지 않는다.
\end{choices}
\end{question}
\answer{2}
